

\drghosh{\textbf{Discussion section.} Empty. Use this to position the finding. Key points: (a) Contrast with the 'flight' literature (Bohra-Mishra \& Massey 2011; Engel \& Ibanez 2007 on Colombia) — they find conflict pushes migration in the SHORT RUN. Your finding is that childhood exposure REDUCES migration in the LONG RUN. The two are reconcilable: people who flee during conflict are economically active adults; the children who stay behind grow up poorer / less educated / less networked, and never migrate. (b) External validity — Nepal's conflict was relatively low-intensity (\textasciitilde{}17,000 deaths) and ended cleanly with a peace accord. Different from Syria, DRC, etc. Be honest about scope. (c) Policy — the conventional wisdom is that post-conflict labor migration eases reintegration and brings remittances. Your finding suggests a missing-migrant problem: the cohorts who would normally have driven Nepal's migration boom are MISSING from migration flows. This has implications for remittance flows, household welfare, and the aggregate labor market.}

